% Candidate insertion only. This file has NOT been inserted into main.tex.
% Confirm that the first 987/971 archived records correspond to the paper's
% reported Monte Carlo batch before using the exact percentile values.

\emph{Failure modes and delayed cooperative engagements.}
To expose the operating limits of the framework, strict interception failures
are distinguished from delayed but ultimately successful cooperative
engagements. Among the 1,000 trials, 13 trials in Case~1 and 29 trials in
Case~2 do not satisfy the strict completion condition that all assigned
interceptors hit their designated targets within the simulation horizon. For
the successful trials, the median, 95th-percentile, and maximum values of
$E_{co\text{-}time}$ are $0.0492$, $0.0708$, and $0.0900$~s in Case~1,
compared with $0.0767$, $0.1521$, and $0.3167$~s in Case~2. In the upper 5\%
timing-error tail of Case~2, the mean engagement duration increases from
$33.244$ to $36.655$~s, whereas the mean miss distance increases only from
$3.082$ to $3.181$~m. The Spearman correlation between
$E_{co\text{-}time}$ and $E_t$ is $0.426$ in Case~2 but only $0.162$ in
Case~1. Thus, continuous target weaving mainly expands the tail of the
closing-time distribution and reduces the simultaneous-impact margin, rather
than causing a broad loss of terminal accuracy. This limitation motivates
time-to-go dispersion monitoring, online reassignment, and reserve-interceptor
activation in practical deployments.

% Suggested figure caption if failure_delay_analysis.pdf is inserted:
% \caption{Upper-tail analysis of delayed cooperative engagement: (a) empirical
% exceedance probability of the temporal-coordination error and (b) its
% relationship with engagement duration. The 0.10-s line is used only as a
% diagnostic reference and not as an interception-success threshold.}
